% PAGES: 109-110

% This file continues Section 3.2 "Arbitrage-free markets" and the proof of
% Theorem 3.1, both begun before PDF page 109 in another agent's file. The
% paragraph below is the tail end of the case-(II)-arbitrage argument; do
% NOT reopen \begin{proof} here -- it closes the proof already in progress.

An arbitrage of type (II) occurs if and only if there exists an $m \times 1$
positions vector $\Theta$ such that
\begin{equation}
V_{t_0} \;=\; \Theta^t S_{t_0} \;<\; 0 \quad \text{and} \quad
V_\tau^{\omega^k} \;=\; \Theta^t S_\tau^{\omega^k} \;\geq\; 0, \quad \forall\, k = 1:n.
\end{equation}
Note that, since $S_{t_0} = M_\tau Q$, formula (3.15) holds, i.e.,
\[
V_{t_0} \;=\; \sum_{k=1}^n Q_k V_\tau^{\omega^k}.
\]
Then, since $Q_k > 0$ and $V_\tau^{\omega^k} \geq 0$ for all $k = 1:n$, we
find that $V_{t_0} \geq 0$, which contradicts the assumption that
$V_{t_0} < 0$ from the no--arbitrage condition (3.16).
\end{proof}

\medskip
\noindent\rule{\linewidth}{0.4pt}
\medskip

\noindent\textbf{\textit{One Period Binomial Model Example (Continued):}}\\
The one period binomial model is arbitrage--free if and only if
\begin{equation}
d \;<\; e^{r\delta t} \;<\; u.
\end{equation}

To see this, note that, from Theorem 3.1, it follows that the one period
binomial model is arbitrage--free if and only if there exist
$Q = \begin{pmatrix} Q_1 \\ Q_2 \end{pmatrix}$ with $Q_1 > 0$ and $Q_2 > 0$
such that
\begin{equation}
M_\tau Q \;=\; S_{t_0},
\end{equation}
where $S_{t_0}$ is the price vector of the securities at time $t_0$ given by
(3.4) and $M_\tau$ is the payoff matrix at time $\tau$ given by (3.5).

The equality (3.18) can be written as follows:
\[
\begin{pmatrix} e^{r\delta t} & e^{r\delta t} \\ uS_0 & dS_0 \end{pmatrix}
\begin{pmatrix} Q_1 \\ Q_2 \end{pmatrix}
\;=\;
\begin{pmatrix} 1 \\ S_0 \end{pmatrix}
\quad\Longleftrightarrow\quad
\begin{cases}
e^{r\delta t} Q_1 \;+\; e^{r\delta t} Q_2 \;=\; 1 \\
uS_0 Q_1 \;+\; dS_0 Q_2 \;=\; S_0
\end{cases}
\]
which can be written as
\begin{align}
Q_1 + Q_2 &= e^{-r\delta t}; \\
uQ_1 + dQ_2 &= 1.
\end{align}

By multiplying (3.19) by $d$ and subtracting the resulting equation from
(3.20) and solving for $Q_1$, we obtain that
\begin{equation}
Q_1 \;=\; \frac{1 - de^{-r\delta t}}{u - d} \;=\; e^{-r\delta t}\,\frac{e^{r\delta t} - d}{u - d}.
\end{equation}

From (3.19) and (3.21), we find that
\begin{equation}
Q_2 \;=\; \frac{ue^{-r\delta t} - 1}{u - d} \;=\; e^{-r\delta t}\,\frac{u - e^{r\delta t}}{u - d}.
\end{equation}

Recall that $d < u$ and therefore $u - d > 0$. Then, from (3.21) and (3.22),
if follows that $Q_1 > 0$ and $Q_2 > 0$ if and only if $d < e^{r\delta t} <
u$. From Theorem 3.1, we conclude that the one period binomial model is
arbitrage--free if and only if $d < e^{r\delta t} < u$, which is what we
wanted to show; cf. (3.17).

The intuition behind the no--arbitrage condition (3.17) is as follows: if
(3.17) were not satisfied, then the following arbitrage opportunities would
occur:

\begin{itemize}
\item If $d \geq e^{r\delta t}$, then the asset will appreciate at least at
the risk free rate even if its value at time $\tau$ is $dS_0$, the lower of
the two possible values for the asset. The arbitrage follows from a
``buy low, sell high'' strategy: borrowing cash and purchasing the asset is
guaranteed not to lose money in either state of the market at time $\tau$,
and will have a positive payoff if the value of the asset at time $\tau$ is
$uS_0$. This is an arbitrage opportunity of type (I).

\item If $e^{r\delta t} \geq u$, then the asset will appreciate at most at
the risk free rate even if its value at time $\tau$ is $uS_0$, the higher of
the two possible values for the asset. The arbitrage follows from a
``buy low, sell high'' strategy: shorting the asset and investing the cash
at the risk--free rate is guaranteed to not lose money in either state of
the market at time $\tau$, and will have a positive payoff if the value of
the asset at time $\tau$ is $dS_0$. This is an arbitrage opportunity of type
(I). \hfill$\square$
\end{itemize}

\medskip
\noindent\rule{\linewidth}{0.4pt}
\medskip

The result below is a special case of the no--arbitrage Law of One Price for
arbitrage--free one period market models; see Stefanica [36, 37] for a
general form of the Law of One Price and more applications.

\begin{theorem}
(The Law of One Price) Consider an arbitrage--free one period market model
with $m$ securities and $n$ market states. If two portfolios have the same
value in every state of the market at time $\tau > t_0$, then they must
have the same value at time $t_0$.
\end{theorem}

\begin{proof}
Let $V_{t_0}$ and $W_{t_0}$ be the values at time $t_0$ of two portfolios
which have the same value in every state of the market at time $\tau$, i.e.
$V_\tau^{\omega^k} = W_\tau^{\omega^k}$, for all $1 \leq k \leq n$. We will
show that if the one period market model is arbitrage--free, then $V_{t_0} =
W_{t_0}$.

If $V_{t_0} \neq W_{t_0}$, assume, for example, that $V_{t_0} < W_{t_0}$ and
consider a portfolio with a long position in the first portfolio and a short
position in the second portfolio. The value at time $t_0$ of this portfolio
is $V_{t_0} - W_{t_0} < 0$, while its value in any state $\omega^k$ at time
$\tau$ is $0$, since $V_\tau^{\omega^k} - W_\tau^{\omega^k} = 0$ for all $k =
1:n$. This is an arbitrage of type (II), which contradicts the fact that the
one period market model was arbitrage--free. The contradiction comes from
the assumption that $V_{t_0} \neq W_{t_0}$, and therefore we conclude that
$V_{t_0} = W_{t_0}$.
\end{proof}

The value of a replicable derivative security in an arbitrage--free market
can be determined as follows:

\begin{lemma}
Consider an arbitrage--free one period market model with $m$ securities and
$n$ market states. Let $Q_k > 0$, $k = 1:n$, be the state prices, and let $Q
= (Q_k)_{k=1:n}$. Let $s_\tau$ be the $1 \times n$ price vector at time
$\tau$ of a replicable security. The value $s_{t_0}$ at time $t_0$ of the
replicable security is
\begin{equation}
s_{t_0} \;=\; s_\tau Q \;=\; \sum_{k=1}^n Q_k s_\tau^{\omega^k},
\end{equation}
where $s_\tau^{\omega^k}$ is the value of the replicable security at time
$\tau$ if state $\omega^k$ occurs.
\end{lemma}

% Proof of Lemma 3.3 continues on PDF page 111, covered by sec03_c2_2.tex.
